\section{Constante de Euler--Kronecker del car\'acter dodecaf\'asico}

La factorizaci\'on de Dedekind

\begin{equation}
\zeta_{\mathbb Q(\sqrt3)}(s)=\zeta(s)L(s,\chi_{12})
\label{eq:dedekind-sqrt3}
\end{equation}

da, al comparar los t\'erminos constantes en \(s=1\),

\begin{equation}
\gamma_{\mathbb Q(\sqrt3)}
=\gamma+\frac{L'(1,\chi_{12})}{L(1,\chi_{12})}
=1.0539656082484825800\ldots.
\label{eq:euler-kronecker}
\end{equation}

Esta constante representa el t\'ermino finito normalizado en el polo de la
funci\'on zeta del campo real dodecaf\'asico. No constituye una variante de
\(\gamma\) de Euler, sino la correcci\'on aritm\'etica inducida por
\(\chi_{12}\).

En efecto, escribiendo \(u=s-1\),
\[
\zeta(s)=u^{-1}+\gamma+O(u),
\qquad
L(s,\chi_{12})=L_1+L_1'u+O(u^2),
\]
se obtiene
\begin{equation}
\zeta_{\mathbb Q(\sqrt3)}(s)
=\frac{L_1}{u}+\bigl(\gamma L_1+L_1'\bigr)+O(u).
\label{eq:ek-laurent-derivation}
\end{equation}
La constante de Euler--Kronecker es el cociente entre el t\'ermino finito y
el residuo; dividir \cref{eq:ek-laurent-derivation} por \(L_1\) prueba
\cref{eq:euler-kronecker}. Esta derivaci\'on explica simult\'aneamente el
residuo
\[
\operatorname*{Res}_{s=1}\zeta_{\mathbb Q(\sqrt3)}(s)
=\frac{\log(2+\sqrt3)}{\sqrt3}
\]
y la correcci\'on finita. El valor num\'erico se certifica evaluando
\(L(1,\chi_{12})\) y \(L'(1,\chi_{12})\) con la misma tabla residual y una
cota de cola; no se importa de una tabla de campos cuadr\'aticos.

\section{Constantes aritm\'eticas asociadas a los n\'umeros primos}

Una vez seleccionados los ciclos aritm\'eticos irreducibles, se obtienen dos
familias distintas:

\begin{align}
B_1&=\lim_{x\to\infty}\left(
\sum_{p\le x}\frac1p-\log\log x
\right),
\label{eq:meissel-mertens}\\
C_2&=\prod_{p>2}\frac{p(p-2)}{(p-1)^2}.
\label{eq:twin-prime-constant}
\end{align}

\(B_1\) regulariza una suma local; \(C_2\) es una densidad multiplicativa de pares gemelos. Ambos requieren control de cola. No se deducen el uno del otro por compartir los primos.

